\documentclass[aspectratio=169,11pt]{beamer}
% Shared Beamer style for practice contest solution walkthroughs.
\usepackage[utf8]{inputenc}
\usepackage[T1]{fontenc}
\usepackage{lmodern}
\usepackage{amsmath,amssymb}
\usepackage{xcolor}
\usepackage{hyperref}
\usepackage{listings}

\definecolor{CPNavy}{HTML}{172033}
\definecolor{CPBlue}{HTML}{2563EB}
\definecolor{CPGreen}{HTML}{16A34A}
\definecolor{CPOrange}{HTML}{F97316}
\definecolor{CPGray}{HTML}{64748B}
\definecolor{CPLight}{HTML}{F8FAFC}
\definecolor{CPCodeBg}{HTML}{F1F5F9}
\definecolor{CPCodeRule}{HTML}{CBD5E1}

\usetheme{default}
\usefonttheme{professionalfonts}
\setbeamertemplate{navigation symbols}{}
\setbeamersize{text margin left=0.65cm,text margin right=0.65cm}
\setbeamertemplate{itemize item}{\color{CPBlue}$\blacktriangleright$}
\setbeamertemplate{itemize subitem}{\color{CPGray}$\triangleright$}

\setbeamercolor{normal text}{fg=CPNavy,bg=white}
\setbeamercolor{frametitle}{fg=white,bg=CPNavy}
\setbeamercolor{title}{fg=white,bg=CPNavy}
\setbeamercolor{subtitle}{fg=CPLight,bg=CPNavy}
\hypersetup{colorlinks=true,urlcolor=CPBlue}

\setbeamertemplate{frametitle}{%
  \nointerlineskip%
  \begin{beamercolorbox}[wd=\paperwidth,ht=0.88cm,dp=0.22cm,leftskip=0.55cm,rightskip=0.35cm]{frametitle}
    \usebeamerfont{frametitle}\insertframetitle\hfill\small\insertframenumber/\inserttotalframenumber
  \end{beamercolorbox}%
}

\setbeamertemplate{footline}{%
  \leavevmode%
  \hbox{%
    \begin{beamercolorbox}[wd=.58\paperwidth,ht=2.4ex,dp=1ex,leftskip=.5cm]{author in head/foot}%
      \color{CPGray}\insertshorttitle
    \end{beamercolorbox}%
    \begin{beamercolorbox}[wd=.42\paperwidth,ht=2.4ex,dp=1ex,rightskip=.5cm plus1fil]{date in head/foot}%
      \hfill\color{CPGray}\insertshortauthor
    \end{beamercolorbox}%
  }%
}

\setbeamertemplate{title page}{%
  \vspace*{1.25cm}
  \begin{beamercolorbox}[wd=\paperwidth,sep=0.65cm,leftskip=0.15cm,rightskip=0.15cm]{title}
    {\color{white}\Huge\bfseries\inserttitle\par}
    \vspace{0.35cm}
    {\color{CPLight}\Large\insertsubtitle\par}
    \vspace{0.6cm}
    {\color{CPLight}\large\insertauthor\par}
  \end{beamercolorbox}
  \vfill
}

\newcommand{\sectiontag}[1]{\textbf{\color{CPBlue}#1}}
\newenvironment{tightitemize}{\begin{itemize}\small\setlength{\itemsep}{0.18cm}}{\end{itemize}}
\lstdefinestyle{cpcode}{
  language=Python,
  basicstyle=\ttfamily\tiny,
  keywordstyle=\color{CPBlue}\bfseries,
  commentstyle=\color{CPGray}\itshape,
  stringstyle=\color{CPGreen},
  backgroundcolor=\color{CPCodeBg},
  frame=single,
  framerule=0.25pt,
  rulecolor=\color{CPCodeRule},
  showstringspaces=false,
  columns=fullflexible,
  keepspaces=true,
  breaklines=true,
  breakatwhitespace=false,
  tabsize=4,
  xleftmargin=0.08cm,
  xrightmargin=0.08cm,
  aboveskip=0.08cm,
  belowskip=0.08cm
}


\title[snowflakes]{Snowflakes}
\subtitle{Day 1 solution walkthrough}
\author[Solution Deck]{Competitive Programming Bootcamp}
\date{}

\begin{document}

\begin{frame}[plain]
  \titlepage
\end{frame}

% BEGIN GENERATED STATEMENT INTERPRETATION
\begin{frame}{Interpret the Word Problem}
\small
\sectiontag{Statement excerpts:}
\begin{tightitemize}
  \item \textit{``bags of uniqueness''}
  \item \textit{``longest contiguous segment of unique snowflakes''}
\end{tightitemize}

\vspace{0.18cm}
\sectiontag{Programming interpretation:}
\begin{tightitemize}
  \item A valid package is a contiguous window with no repeated snowflake ID.
  \item When a duplicate appears inside the current window, the left edge must move past the previous copy.
  \item The largest valid window length is the maximum package size.
\end{tightitemize}
\end{frame}
% END GENERATED STATEMENT INTERPRETATION







\begin{frame}{Problem Model}
\small
\sectiontag{Textbook connection:} CP4 2.3 f. Hash Table (map), Harder

\vspace{0.25cm}
\sectiontag{Core technique:} Sliding window with a hash map

\vspace{0.25cm}
\begin{tightitemize}
  \item We need the longest contiguous sequence with no repeated snowflake ID.
  \item The current window [left, right] is always duplicate-free.
  \item A dictionary records the most recent index of each snowflake.
\end{tightitemize}
\end{frame}

\begin{frame}{How to Recognize the Approach}
\begin{tightitemize}
  \item Contiguous + unique elements points to a sliding window.
  \item A hash table tells us whether a duplicate is inside the active window.
  \item Moving left only forward keeps the total work linear.
\end{tightitemize}
\end{frame}

\begin{frame}{Algorithm}
\begin{tightitemize}
  \item Initialize left = 0, best = 0, and an empty last\_seen dictionary.
  \item Scan each snowflake as the right endpoint.
  \item If this snowflake was last seen at or after left, move left past that position.
  \item Update last\_seen for the current snowflake.
  \item Update best with right - left + 1.
\end{tightitemize}
\end{frame}

\begin{frame}{Walkthrough of the Key State}
\begin{tightitemize}
  \item Duplicates before left are irrelevant because they are outside the current window.
  \item A duplicate inside the window forces the shortest safe shift of left.
  \item The window never contains two copies of the same ID after the update.
\end{tightitemize}
\end{frame}

\begin{frame}{Why the Algorithm Is Correct}
\begin{tightitemize}
  \item The invariant is that the maintained window is always duplicate-free.
  \item When a duplicate appears, any valid window ending at right must start after the previous copy.
  \item Since every right endpoint is considered and best stores the largest valid window, the final answer is optimal.
\end{tightitemize}
\end{frame}

\begin{frame}{Complexity and Implementation Notes}
\small
\sectiontag{Complexity:} O(N) expected time per test case and O(N) memory.

\vspace{0.3cm}
\begin{tightitemize}
  \item Use sys.stdin.buffer.read for fast input across many test cases.
  \item Store positions, not counts, because the exact previous index determines the new left boundary.
\end{tightitemize}
\end{frame}



% BEGIN GENERATED CODE WALKTHROUGH
\begin{frame}[fragile]{Code: Bulk Input Setup}
\scriptsize
\sectiontag{Source lines:} \texttt{snowflakes.py:27--37}

\vspace{0.08cm}
\begin{columns}[T,totalwidth=\textwidth]
\begin{column}{0.58\textwidth}
\begin{lstlisting}[style=cpcode]
import sys

def solve():
    data = list(map(int, sys.stdin.buffer.read().split()))
    pos = 0

    test_cases = data[pos]
    pos += 1

    answers = []
\end{lstlisting}
\end{column}
\begin{column}{0.38\textwidth}
\sectiontag{What this chunk does}
\begin{tightitemize}
  \item The solution reads all integers once for fast input.
  \item pos is a cursor into the flat token list.
  \item answers collects one result per test case before printing.
\end{tightitemize}
\end{column}
\end{columns}
\end{frame}

\begin{frame}[fragile]{Code: Start One Window}
\scriptsize
\sectiontag{Source lines:} \texttt{snowflakes.py:39--45}

\vspace{0.08cm}
\begin{columns}[T,totalwidth=\textwidth]
\begin{column}{0.58\textwidth}
\begin{lstlisting}[style=cpcode]
for _ in range(test_cases):
    n = data[pos]
    pos += 1

    last_seen = {}
    left = 0
    best = 0
\end{lstlisting}
\end{column}
\begin{column}{0.38\textwidth}
\sectiontag{What this chunk does}
\begin{tightitemize}
  \item Each test case starts with n snowflake IDs.
  \item last\_seen maps an ID to its most recent index.
  \item left and best describe the current unique window and best answer.
\end{tightitemize}
\end{column}
\end{columns}
\end{frame}

\begin{frame}[fragile]{Code: Slide Past Duplicates}
\scriptsize
\sectiontag{Source lines:} \texttt{snowflakes.py:47--57}

\vspace{0.08cm}
\begin{columns}[T,totalwidth=\textwidth]
\begin{column}{0.58\textwidth}
\begin{lstlisting}[style=cpcode]
for right in range(n):
    snowflake = data[pos]
    pos += 1

    if snowflake in last_seen and last_seen[snowflake] >= left:
        left = last_seen[snowflake] + 1

    last_seen[snowflake] = right

    best = max(best, right - left + 1)
\end{lstlisting}
\end{column}
\begin{column}{0.38\textwidth}
\sectiontag{What this chunk does}
\begin{tightitemize}
  \item right scans each snowflake exactly once.
  \item If a duplicate is still inside the window, left jumps after its old position.
  \item Updating last\_seen after the jump records the newest occurrence.
\end{tightitemize}
\end{column}
\end{columns}
\end{frame}

\begin{frame}[fragile]{Code: Store Case Results}
\scriptsize
\sectiontag{Source lines:} \texttt{snowflakes.py:59--64}

\vspace{0.08cm}
\begin{columns}[T,totalwidth=\textwidth]
\begin{column}{0.58\textwidth}
\begin{lstlisting}[style=cpcode]
        answers.append(str(best))

    print("\n".join(answers))

solve()
\end{lstlisting}
\end{column}
\begin{column}{0.38\textwidth}
\sectiontag{What this chunk does}
\begin{tightitemize}
  \item best is appended after the whole stream for one test case is processed.
  \item A single joined print avoids repeated output calls.
  \item solve() is invoked at the bottom to run the workflow.
\end{tightitemize}
\end{column}
\end{columns}
\end{frame}
% END GENERATED CODE WALKTHROUGH

\begin{frame}{Source}
\small
\sectiontag{Solution file:} \texttt{practice\_contests/day\_1/snowflakes.py}

\vspace{0.3cm}
\sectiontag{Problem:} \url{https://open.kattis.com/problems/snowflakes}

\vspace{0.45cm}
Use this deck with the source code open beside it. The slides explain the reasoning; the code shows the exact input handling and output formatting.
\end{frame}

\end{document}
